\section{The Statutory Argument}
\label{sec:statutory}

\subsection{Two Tasks, Two Methods}

The redistricting enterprise involves two distinct computational tasks:

\begin{enumerate}
\item \textbf{Plan generation}: produce the map that will govern elections
      for the next decade.
\item \textbf{Plan evaluation}: determine whether an enacted map is a
      partisan outlier relative to the space of algorithmically neutral plans.
\end{enumerate}

These tasks have different requirements and different optimal methods.

\paragraph{Plan generation.}
For a statutory plan, the overriding requirements are:
(a) determinism (the same inputs always yield the same output);
(b) optimality (the plan minimises the objective function);
(c) transparency (any party can verify the computation in hours, not days).
The \cs{} algorithm satisfies all three.
It requires no chain convergence certification; the stopping criterion
($T = 600$ consecutive non-improving seeds) is a finite, checkable computation.
MCMC plan generation fails (a) without fixing a random seed, fails (b)
without convergence proof, and fails (c) for large states.

\paragraph{Plan evaluation.}
For auditing an enacted map---whether drawn algorithmically or by a
legislature---the requirements are:
(a) an unbiased sample from the uniform distribution over valid plans;
(b) sufficient sample size for reliable quantile estimates;
(c) transparent stopping rules backed by convergence diagnostics.
\recom{} ensemble sampling with certified convergence satisfies all three.
\cs{} is not applicable here; it generates one plan, not a sample.

\subsection{Certified Convergence as an Evidentiary Standard}

Drawing on the diagnostic thresholds developed in Sections~\ref{sec:rhat}--\ref{sec:mixing-time},
we propose the following minimum standards for an ensemble to be admitted
as expert evidence in redistricting litigation:

\begin{enumerate}
\item $\rhat < 1.1$ for both PP and $D$-seats across $m \geq 5$ parallel chains.
\item $\ess > 500$ for PP and $\ess > 300$ for $D$-seats.
\item $\ham(1) < 0.95$.
\item Chain length $\geq$ the state-specific minimum in Table~\ref{tab:mixing-times}.
\end{enumerate}

Condition (3) is a sanity check on the proposal distribution; if
$\ham(1) \geq 0.95$, the chain is moving too slowly to be a useful
diagnostic tool.

\subsection{Why 10{,}000 Steps is the Practical Minimum}

From Table~\ref{tab:diagnostics}, $\ess > 500$ requires at minimum
10{,}000 steps for all six study states.
$\rhat < 1.1$ is achieved at 10{,}000 steps for four of six states;
Texas and California require more.

We therefore recommend 10{,}000 steps as the minimum for states with
$k \leq 20$, with the step count scaling as $500 k$ for $k > 20$.
Under this formula:
\[
  n_{\rm min}(k) = \begin{cases} 10{,}000 & k \leq 20 \\ 500k & k > 20. \end{cases}
\]

For Texas ($k = 38$): $n_{\rm min} = 19{,}000$.
For California ($k = 52$): $n_{\rm min} = 26{,}000$.
These are consistent with Table~\ref{tab:mixing-times}.

\subsection{ConvergenceSweep is Cheaper and More Transparent}

For the specific task of generating the DIA statutory map,
the \cs{} algorithm (B.16) is unambiguously preferable:
\begin{itemize}
\item Runtime: $O(T \cdot n)$ with $T = 600$ and $n$ the number of tracts.
      For Texas: $\approx 600 \times 8{,}000 = 4.8M$ METIS calls.
      In practice, \cs{} completes in under 60 seconds for any state.
\item Certification: the stopping criterion is a finite inequality check,
      not an asymptotic convergence argument.
\item Auditability: the sequence of seeds, the objective values, and the
      final map are all deterministic given the census release identifier.
\end{itemize}

Ensemble sampling requires hours of computation and a convergence
certificate based on asymptotic theory.
For plan generation, \cs{} is both cheaper and more transparent.
